% ============================================================
%  pandoc_header_techbook.tex
%  Technical Book Template — pandoc + lualatex + ltjsreport
%  Updated cover page, header (logo), and footer (copyright)
% ============================================================

%% ─── 1. Color Palette ───────────────────────────────────────
\usepackage{xcolor}
\definecolor{TBNavy}{RGB}{22, 58, 106}
\definecolor{TBAccent}{RGB}{50, 110, 175}
\definecolor{TBLightBlue}{RGB}{220, 232, 248}
\definecolor{TBPalest}{RGB}{236, 243, 252}
\definecolor{TBDarkGray}{RGB}{55, 55, 55}
\definecolor{TBMidGray}{RGB}{120, 120, 120}
\definecolor{TBRule}{RGB}{160, 190, 220}
\definecolor{TBCodeBg}{RGB}{244, 247, 251}
\definecolor{TBLink}{RGB}{0, 70, 140}     % URL・外部リンク（明るい青）
\definecolor{TBInternalLink}{RGB}{50, 110, 175} % 内部リンク（目次・相互参照）

%% ─── 2. フォント設定（ゴシック体） ──────────────────────────
\usepackage{luatexja-fontspec}
% macOS: Hiragino / Linux (CI): Noto Sans CJK JP
\IfFontExistsTF{Hiragino Kaku Gothic ProN W3}{%
  \setmainjfont[BoldFont = Hiragino Kaku Gothic ProN W6]{Hiragino Kaku Gothic ProN W3}%
}{%
  \setmainjfont{Noto Sans CJK JP}%
}
\IfFontExistsTF{TeX Gyre Heros}{%
  \setmainfont[BoldFont = TeX Gyre Heros Bold, ItalicFont = TeX Gyre Heros Italic, Scale = MatchLowercase]{TeX Gyre Heros}%
}{%
  \setmainfont[Scale = MatchLowercase]{Latin Modern Sans}%
}

%% ─── 3. Hyperref 設定 ───────────────────────────────────────
%  リンク色は Makefile の -V colorlinks=true / -V linkcolor=... で制御。
%  pandoc が colorlinks=true を受け取ると hidelinks を出力しなくなり、
%  代わりに linkcolor / urlcolor / citecolor 変数の値をそのまま使う。
%  ここでは追加設定のみ行う。
\AtBeginDocument{%
  \hypersetup{pdfborder={0 0 0}}%
}

%% ─── 3b. ドキュメント初期設定 ──────────────────────────────────
%  \TBCopyYear は year_meta.lua が header-includes に注入する。
\AtBeginDocument{%
  \setlength{\footskip}{10mm}%
  \addtocontents{toc}{\protect\markboth{目次}{目次}}%
}

%% ─── 4. タイトルページ ───────────────────────────────────────
%  右上にロゴ、タイトル + 横罫線 + 版数 + 日付、左下に copyright。
%  \docversion は version_meta.lua が header-includes に注入する。
\makeatletter
\renewcommand{\maketitle}{%
  \newpage
  \thispagestyle{empty}%
  % ロゴ（右上）— 色付き帯 + "MB-POST" テキスト
  \noindent\hfill
  \origincludegraphics[height=18mm]{logo.pdf}%
  \hspace{6pt}%
  \raisebox{7mm}{\Large\bfseries\color{TBDarkGray}MB-POST}%
  \par
  \vspace*{5cm}
  % タイトル
  \noindent{\fontsize{26}{36}\selectfont\bfseries\@title\par}%
  \vspace{2pt}
  \noindent{\color{TBNavy}\rule{\linewidth}{1.5pt}\par}%
  \vspace{10pt}
  % 版数・日付
  \noindent{\large\color{TBDarkGray}\docversion\par}%
  \vspace{4pt}
  \noindent{\large\color{TBDarkGray}\@date\par}%
  \vfill
  % Copyright（左下）
  \noindent{\small\itshape\color{TBMidGray}%
    Copyright \textcopyright\ \TBCopyYear\ \@author\par}%
  \newpage
}
\makeatother

%% ─── 5. 章・節タイトル ───────────────────────────────────────
%  [explicit] で見出しテキスト #1 をフォーマット内で参照可能にする
\usepackage[explicit]{titlesec}

%% ── 5a. Chapter（番号付き）: ネイビー帯 + 薄青タイトル背景
\titleformat{\chapter}[block]
  {\normalfont\bfseries}
  {}
  {0pt}
  {%
    \vbox{%
      \hbox to \linewidth{%
        \colorbox{TBNavy}{\parbox{\dimexpr\linewidth-2\fboxsep\relax}{%
          \vspace{5pt}%
          \hspace{10pt}%
          {\color{white}\large\mdseries 第\thechapter 章}%
          \vspace{5pt}%
        }}%
      }%
      \nointerlineskip
      \hbox to \linewidth{%
        \colorbox{TBLightBlue}{\parbox{\dimexpr\linewidth-2\fboxsep\relax}{%
          \vspace{12pt}%
          \hspace{14pt}%
          {\color{TBNavy}\LARGE\bfseries #1}%
          \vspace{14pt}%
        }}%
      }%
    }%
  }
  [\thispagestyle{fancy}]

\titlespacing*{\chapter}{0pt}{0pt}{24pt}

%% ── 5b. Chapter*（目次など番号なし章）: 薄青タイトル背景のみ
\titleformat{name=\chapter,numberless}[block]
  {\normalfont\bfseries}
  {}
  {0pt}
  {%
    \markboth{#1}{#1}%
    \colorbox{TBLightBlue}{\parbox{\dimexpr\linewidth-2\fboxsep\relax}{%
      \vspace{12pt}%
      \hspace{14pt}%
      {\color{TBNavy}\LARGE\bfseries #1}%
      \vspace{14pt}%
    }}%
  }
  [\thispagestyle{fancy}]

\titlespacing*{name=\chapter,numberless}{0pt}{0pt}{24pt}

%% ── 5c. Section: 薄青背景 + 青左縦線
\titleformat{\section}[block]
  {\normalfont\large\bfseries}
  {}
  {0pt}
  {%
    \colorbox{TBLightBlue}{\parbox{\dimexpr\linewidth-2\fboxsep\relax}{%
      \vspace{4pt}%
      {\color{TBAccent}\rule[-3pt]{4pt}{16pt}}%
      \hspace{8pt}%
      {\color{TBNavy}\thesection\hspace{8pt}#1}%
      \vspace{5pt}%
    }}%
  }

\titlespacing*{\section}{0pt}{18pt}{8pt}

%% ── 5d. Subsection: 最薄青背景 + 細左縦線
\titleformat{\subsection}[block]
  {\normalfont\normalsize\bfseries}
  {}
  {0pt}
  {%
    \colorbox{TBPalest}{\parbox{\dimexpr\linewidth-2\fboxsep\relax}{%
      \vspace{3pt}%
      {\color{TBAccent}\rule[-2pt]{3pt}{11pt}}%
      \hspace{7pt}%
      {\color{TBNavy}\thesubsection\hspace{6pt}#1}%
      \vspace{4pt}%
    }}%
  }

\titlespacing*{\subsection}{0pt}{12pt}{5pt}

%% ── 5e. Subsubsection: 青縦線のみ（背景なし）
\titleformat{\subsubsection}[hang]
  {\normalfont\normalsize\bfseries\color{TBNavy}}
  {{\color{TBAccent}\rule[-2pt]{2pt}{10pt}}\hspace{6pt}\thesubsubsection}
  {0.6em}
  {#1}

\titlespacing*{\subsubsection}{0pt}{10pt}{3pt}

%% ─── 6. ヘッダー・フッター ───────────────────────────────────
%  ヘッダー左: 節番号 + 節タイトル（章の先頭ページは章タイトル）
%  ヘッダー右: MB-POST ロゴ
%  フッター左: Copyright © YEAR AUTHOR
%  フッター右: | ページ番号

\usepackage{fancyhdr}
\setlength{\headheight}{14mm}

\makeatletter

%  ltjsreport の \chaptermark / \sectionmark を上書き
\renewcommand{\chaptermark}[1]{%
  \ifnum\value{chapter}>0
    \markboth{\thechapter\enspace #1}{\thechapter\enspace #1}%
  \else
    \markboth{#1}{#1}%
  \fi
}
\renewcommand{\sectionmark}[1]{%
  \markright{\thesection\enspace #1}%
}

\pagestyle{fancy}
\fancyhf{}

%  左: 節番号 + 節タイトル（\rightmark）
\fancyhead[LE,LO]{\small\color{TBMidGray}\nouppercase{\rightmark}}
%  右: ロゴ + "MB-POST" テキスト
\fancyhead[RE,RO]{%
  \origincludegraphics[height=10mm]{logo.pdf}%
  \hspace{5pt}%
  \raisebox{3.5mm}{\small\bfseries\color{TBDarkGray}MB-POST}%
}

\renewcommand{\headrulewidth}{0.5pt}
\renewcommand{\headrule}{%
  {\color{TBRule}\hrule width\headwidth height\headrulewidth}%
  \vskip-\headrulewidth
}
\renewcommand{\footrulewidth}{0pt}

%  フッター左: copyright
\fancyfoot[LE,LO]{\small\itshape\color{TBMidGray}%
  Copyright \textcopyright\ \TBCopyYear\ \@author}
%  フッター右: | ページ番号
\fancyfoot[RE,RO]{\small\color{TBMidGray}\textbar\enspace\thepage}

%  章の先頭ページ（plain スタイル）: 左に章タイトル、右にロゴ、フッターは同じ
\fancypagestyle{plain}{%
  \fancyhf{}%
  \fancyhead[LE,LO]{\small\color{TBMidGray}\nouppercase{\leftmark}}%
  \fancyhead[RE,RO]{%
    \origincludegraphics[height=10mm]{logo.pdf}%
    \hspace{5pt}%
    \raisebox{3.5mm}{\small\bfseries\color{TBDarkGray}MB-POST}%
  }%
  \fancyfoot[LE,LO]{\small\itshape\color{TBMidGray}%
    Copyright \textcopyright\ \TBCopyYear\ \@author}%
  \fancyfoot[RE,RO]{\small\color{TBMidGray}\textbar\enspace\thepage}%
  \renewcommand{\headrulewidth}{0.5pt}%
  \renewcommand{\headrule}{%
    {\color{TBRule}\hrule width\headwidth height\headrulewidth}%
    \vskip-\headrulewidth}%
}

\makeatother

%% ─── 7. 表の行間正規化・セル内改行 ─────────────────────────────
%  longtable ヘッダの minipage に linestretch が適用されて
%  行間が異常に広がる問題を防ぐ
\usepackage{makecell}
\usepackage{colortbl}
\colorlet{TBTableAlt}{TBPalest}
\usepackage{etoolbox}
\AtBeginEnvironment{longtable}{%
  \setstretch{1.1}%
  \small%
}
\AtBeginEnvironment{tabular}{%
  \setstretch{1.1}%
  \small%
}

%% ─── 7b. 章の先頭ページへの fancy スタイル強制（patchcmd） ─────
%  ltjsreport の \chapter は \thispagestyle{plain} を呼び出す。
%  titleformat の after-code と二重に対策することで確実に適用する。
\makeatletter
\patchcmd{\chapter}{\thispagestyle{plain}}{\thispagestyle{fancy}}{}{}
\makeatother

%% ─── 8. コードブロック ───────────────────────────────────────
\usepackage{listings}
\lstset{
  basicstyle        = \small\ttfamily\color{TBDarkGray},
  backgroundcolor   = \color{TBCodeBg},
  frame             = single,
  framerule         = 0.5pt,
  rulecolor         = \color{TBRule},
  xleftmargin       = 8pt,
  xrightmargin      = 8pt,
  framexleftmargin  = 8pt,
  framexrightmargin = 8pt,
  breaklines        = true,
  breakatwhitespace = false,
  keepspaces        = true,
  columns           = flexible,
  showstringspaces  = false,
  keywordstyle      = \bfseries\color{TBNavy},
  commentstyle      = \itshape\color{TBMidGray},
  stringstyle       = \color{TBAccent},
  numbers           = none,
  aboveskip         = 8pt,
  belowskip         = 8pt,
}

%% ─── 9. 図キャプション ───────────────────────────────────────
\usepackage{caption}
\captionsetup{
  font          = small,
  labelfont     = {bf, color=TBDarkGray},
  labelsep      = period,
  textfont      = {color=TBDarkGray},
  justification = centering,
  skip          = 6pt,
}

%% ─── 10. リスト余白調整 ──────────────────────────────────────
%  linestretch=1.2 環境（1行≒14.4pt）に合わせた設定。
%  pandoc デフォルトの \parskip=6pt はリスト前の空白を広く見せるため 3pt に抑制。
%  topsep はリスト開始・終了時の追加余白（\parskip と合算されるため小さめに）。
\usepackage{enumitem}
\setlength{\parskip}{3pt}   % pandoc デフォルト 6pt → 3pt（リスト前の過剰空白を抑制）
\setlist{topsep=2pt, itemsep=2pt, parsep=0pt, partopsep=0pt}
\setlist[itemize,1]{leftmargin=*}
\setlist[enumerate,1]{leftmargin=*}

%% ─── 11. URL スタイル ────────────────────────────────────────
\usepackage{url}
\urlstyle{same}

%% ─── 12. 図のボーダー ────────────────────────────────────────
%  \fboxsep / \fboxrule をグループ内で変更し、グローバルへの
%  漏れ出しを防ぐ（colorbox の幅計算ずれを防止）
\makeatletter
\def\maxwidth{\ifdim\Gin@nat@width>\linewidth\linewidth\else\Gin@nat@width\fi}
\makeatother

\let\origincludegraphics\includegraphics
\renewcommand{\includegraphics}[2][]{%
  \begingroup
    \setlength{\fboxsep}{5pt}%
    \setlength{\fboxrule}{2pt}%
    \fcolorbox{TBRule}{white}{%
      \origincludegraphics[%
        width=\dimexpr\linewidth-2\fboxsep-2\fboxrule\relax,%
        height=0.82\textheight,keepaspectratio%
      ]{#2}%
    }%
  \endgroup
}

%% ─── 13. 長い英単語・URLの折り返し改善 ──────────────────────
\sloppy
\emergencystretch=3em

%% ─── 14. 図の配置: Markdown 記述位置に強制固定 ──────────────
%  pandoc は figure 環境を [htbp]（浮動）で生成する。
%  float パッケージの H 指定（HERE, exactly）で上書きし、
%  ソース中の挿入位置通りに画像を配置する。
%
%  [H] 配置では \intextsep / \textfloatsep が適用されないため、
%  \AtBeginEnvironment / \AtEndEnvironment で figure 前後に
%  明示的な余白を追加する。
\usepackage{float}
\floatplacement{figure}{H}

%  figure 環境の前後に余白を追加（[H] 配置専用の対処）
%  \addvspace は直前に同種の空白があれば加算しない（二重余白防止）
\AtBeginEnvironment{figure}{\addvspace{10pt}}
\AtEndEnvironment{figure}{\addvspace{8pt}}
